
\mysection{Overview}

\tool is a memory-efficient GPU sharing system for cost-effective multi-LLM serving.
Figure~\ref{fig:design-overview} shows its system architecture and design overview. \tool receives inference requests at its frontend, which routes them to the corresponding LLMs for processing. 
To improve cost efficiency, \tool serves LLMs with flexible combinations of space and time sharing. For instance, high-rate models may occupy GPUs exclusively, while multiple low-rate models can be colocated on a single GPU or a model-parallel GPU group. In \tool, a GPU group represents the strict scheduling boundary\textemdash a set of GPUs tightly coupled to jointly serve one model instance. \tool's global scheduler treats each of these groups as a distinct, indivisible resource unit. Idle models are temporarily evicted to CPU DRAM and reactivated when new requests arrive. \tool continuously adjusts its sharing strategies based on runtime workload to maximize overall SLO attainment. 

\tool achieves this with two key designs for flexible, demand-aware cross-model memory sharing and coordination. First, it introduces a balloon driver, referred to as \codeIn{kvcached} (\S\ref{sec:sys}) that enables flexible and efficient memory sharing between models, forming the foundation for rapid adaptation to workload variations and policy changes.
Second, \tool optimizes the overall resource allocation through a memory-centric control plane that uses a combination of (1) load-aware model placement (\S\ref{sec:sch-global}) that maximizes memory headroom for GPUs to facilitate ballooning and (2) slack-aware request arbitration that prioritizes requests on each GPU for maximizing SLO attainment (\S\ref{sec:sch-gpu}).

\MyPara{Relation with autoscaling.}
\tool is orthogonal to and can be seamlessly integrated with autoscaling~\cite{blazescale, sgl‐project-ome, ai-dynamo_2025, llm-d_2025}.
\tool focuses on efficiently sharing GPUs between models (replicas) while autoscaling automatically scales up/down the number of model replicas.

\begin{figure}
    \centering
    \includegraphics[width=0.9\linewidth]{figs/overview.pdf}
    \caption{The system architecture and design overview of \tool.}
    \vspace{-0mm}
    \label{fig:design-overview}
\end{figure}

